% concatenated from flexible.tex
\documentclass[a4paper]{amsart}
\usepackage{amsmath, amsfonts,amsthm,amssymb,amscd, verbatim,graphicx,color,multirow,booktabs, caption, bm,chngcntr,adjustbox,longtable,mathtools,microtype, xfrac}
\usepackage[top = 78pt, bottom = 72pt, inner=108pt, outer=108pt]{geometry}% the inner and outer shall be reduced to 72pt on the final version
\usepackage{hyperref}
\usepackage[dvipsnames, table]{xcolor}
\definecolor{LinkBlueSlate}{HTML}{2F4D73}
\hypersetup{colorlinks=true, linkcolor=Maroon, citecolor=OliveGreen, urlcolor=LinkBlueSlate}
\usepackage{tikz, tikz-cd, quiver}

\renewcommand\qedsymbol{%
	\color{OliveGreen} \vrule height5pt depth0pt width5pt}

\usepackage[capitalise,noabbrev]{cleveref}
\usepackage[shortlabels]{enumitem}
\renewcommand{\theequation}{\Alph{equation}}

\baselineskip=13pt % comfortable reading

\usepackage{todonotes}
\newcommand{\authorcomment}[4]{\todo[color=#1,#2]{#3: #4}{}\PackageWarning{todonotes}{#3: #4}}
\newcommand{\marco}[2][]{\authorcomment{red!40}{#1}{Marco}{#2}}
\newcommand{\andoni}[2][]{\authorcomment{blue!40}{#1}{Andoni}{#2}}

\renewcommand{\wr}{ \mathop{\mathrm{wr}} }
\newcommand{\g}{\mathrm{girth}}
\newcommand{\len}{\mathrm{length}}
\newcommand{\Aut}{\mathrm{Aut}}
\newcommand{\B}{\mathbf{B}}
\newcommand{\Z}{\mathbb{Z}}
\newcommand{\F}{\mathbb{F}}
\newcommand{\X}{\mathbf{X}}
\newcommand{\0}{\mathbf{0}}
\newcommand{\J}{\mathcal{J}}
\newcommand{\M}{\mathcal{M}}

\newtheorem{thm}{Theorem}[]
\newtheorem{lemma}[thm]{Lemma}
\newtheorem{corollary}[thm]{Corollary}
\newtheorem{proposition}[thm]{Proposition}
	
\newtheorem{thmx}{Theorem}[]
\renewcommand{\thethmx}{\Alph{thmx}}
\newtheorem{corx}[thmx]{Corollary}

\theoremstyle{definition}
\newtheorem{de}[thm]{Definition}
\newtheorem{remark}[thm]{Remark}
\newtheorem{construction}[thm]{Construction}
\newtheorem{question}[thmx]{Question}

\title{Flexible $3$-valent graphs of arbitrary girth}

\author[M.~Barbieri]{Marco Barbieri}
\address{Fakulteta za Matematiko in Fiziko, Univerza v Ljubljani, 1000 Jadranska, 1000 Ljubljana, Slovenia} 
\email{marco.barbieri@fmf.uni-lj.si}

\author[A.~Zozaya]{Andoni Zozaya} 
\address{Department of Statistics, Computer Science and Mathematics, Public University of Navarra (UPNA) \& Institute for Advanced Materials and Mathematics (INAMAT$^2$), \linebreak Arrosadia Campus, 31006 Pamplona, Spain}
\email{andoni.zozaya@unavarra.es}

\subjclass[2020]{primary: 20B25; secondary: 05C12, 20E05.}
\keywords{Flexible, vertex-transitive, $3$-valent, amalgamated product}

\title{Flexible $3$-valent graphs of even girth}

\thanks{\textit{Funding:} The first author is supported by the Slovenian Research Agency grant J1-50001 and research programme P1-0222. He is also a member of the GNSAGA INdAM research group and gratefully acknowledges its support. The second author is supported by Spanish Government, grant PID2020-117281GB-I00 (partly funded with ERDF), and by the Public University of Navarra, project 244 \textit{Álgebra. Aplicaciones}. The work on this paper began while the first author was visiting Prof.\ Primo\v{z} Poto\v{c}nik at IMFM in Ljubljana, and the second author was employed at the University of Ljubljana.}

\begin{document}	
	
    \begin{abstract}
        We prove the existence of a connected flexible $3$-valent vertex-transitive graph of girth $2\ell$ for every integer $\ell$. We also give a constructive proof if $\ell$ is prime.
    \end{abstract}
	
    \maketitle

\section{Introduction}\label{sec:intro}

Vertex-transitive $3$-valent graphs fall into three classes according to the number of orbits that their automorphism group induces on the edge set: \emph{arc-transitive graphs} have a single edge orbit, \emph{Cayley graphs} have three, and the remaining case has two. The former two families have been widely studied, while the last family is usually the harshest obstacle while trying to obtain general results for vertex-transitive $3$-valent graphs (for some recent examples of this, we suggest~\cite{BarbieriGrazianSpiga2025,MaSc,   PS21, PSV15, PTV24, Spiga2014}). In this setting, a connected vertex-transitive $3$-valent graph is said to be \emph{flexible} if its automorphism group has exactly two edge-orbits and every vertex-stabiliser has order at least $4$, or, equivalently, each vertex-stabiliser is a noncyclic $2$-group.\footnote{The difficulty in studying \(3\)-valent graphs with two edge orbits lies in the fact that the order of a vertex stabilizer can be arbitrarily large. In contrast, when the stabilizers are cyclic of order \(2\), their behaviour can be completely determined by their action on the neighbourhood of the fixed vertex. This is the reason for the exclusion in the definition of flexible.}

A well-studied parameter in graph theory is the \emph{distinguishing number} of a graph, defined as the minimum number of colours that must be assigned to the vertices of the graph so that no nontrivial automorphism preserves the colouration (see~\cite{AlbertsonCollin}). With only finitely many exceptions, two colours suffice for all $3$-valent vertex-transitive graphs ~\cite{HIKST}. This raises a natural question: \emph{what is the minimum number of vertices which need to be coloured to break all symmetries?} The parameter capturing this idea is called the \emph{$2$-distinguishing cost} of the graph, and it has been introduced in~\cite{Boutin}. It is immediate that the cost is \(1\) in the Cayley case (see~\cite{Sabidussi}), while for arc-transitive cubic graphs it is bounded between \(2\) and \(5\) (see~\cite{Tutte}).  The flexible case has only recently begun to be analysed~\cite{Ademir,ILTW}, and current evidence suggests that, for flexible cubic graphs, the $2$-distinguishing cost is governed by the girth. This led them to pose the following~\cite[Question~3.6]{Ademir}: \emph{does there exist a connected flexible $3$-valent vertex-transitive graph of girth $g$ for every positive integer $g$?} We answer this question in the affirmative for all even values of \(g\). 

\begin{thmx}\label{thm}
    For every integer $\ell$, there exists a finite connected vertex-transitive $3$-valent and flexible graph $\Gamma_\ell$ whose girth is $2\ell$.
\end{thmx}

The analogous existence problems in the one- and three-edge-orbit families have been addressed in~\cite{arc-transitive} and in~\cite{JS11}, respectively. In contrast, examples of such flexible graphs are elusive: $10$ was the largest girth for which a flexible graph was known. In fact, all flexible graphs of girth at most $6$ are known, as part of the wider classification of all $3$-valent vertex-transitive graphs of girth at most $6$, obtained in~\cite{EibenJajcaySparl, VidaliPotocnik}. Further examples of girth up to $10$ can be found in the census of $3$-valent vertex-transitive graphs built in~\cite{PotocnikSpigaVerret2013}. (To find them, we used the \textsf{GAP} package \texttt{GraphSym}~\cite{Rhys}.\footnote{For girth $7$, the smallest flexible graph has $98$ vertices and database number $13$. For girth $8$, there are two examples on $64$ vertices, with database numbers $4$ and $6$. For girth $9$, there are three graphs on $486$ vertices, whose database numbers are $56$, $59$, and $71$. Finally, for girth $10$, there is a unique graph on $144$ vertices database number $68$.})

In our proof of \cref{thm}, girth cycles necessarily appear with even length (see the proof of \cref{lemma:finiteIndexNormal}). As a consequence, our argument does not extend to the case of odd girths. This naturally leads us to pose the following question.
\begin{question}
    Can we construct an explicit example of a finite connected vertex-transitive $3$-valent and flexible graph whose girth is odd and exceeds $11$?
\end{question}
    
Furthermore, our proof of \cref{thm} is not in general constructive, and thus we have no control on the number of vertices of $\Gamma_\ell$. However, can adapt the proof to make it constructive if $\ell$ is a prime, thus obtaining an upper bound on the number of vertices of $\Gamma_\ell$. Our method relays on finding an explicit finite index subgroup of appropriate girth\footnote{The girth of a normal subgroup $N \unlhd F$ is the girth of the Cayley graph of $F/N$ with respect to the canonical generators.} of the free group. (The study of girth of subgroup is an area of independent interest). % -- see, for instance,~\cite{Matucci,MaPu,Rivin}.)

\begin{thmx}
\label{cor: p-power-nilpotent}
    Let $F_k = \langle x_1, \dots, x_k \rangle $ be the free group of rank $k$, let $\ell$ be a prime number, let $\gamma_{\ell+1}(F_k)$ be the $\ell$-th term of the lower central series of $F_k$, and let $F_k^\ell$ be the group generated by $\ell$-powers. Then, for every $w \in F_k^\ell \gamma_{\ell}(F_k)$, $\len(w) \geq \ell$, and equality is attained if, and only if, $w = x_{i}^{\pm \ell}$, for some index $i \in \{ 1, \dots, k \}$.
\end{thmx}
\begin{corx}
\label{prop:bounds}
    Let $\ell$ be a prime number. Then
    \[ |V\Gamma_\ell| \le 16 \cdot N\left(3,2\ell\right) \,,\]
    where $N(3,2\ell)$ is the order of the largest $3$-generated finite group of exponent $2\ell$.
\end{corx}
Observe that $N(3,2\ell)$ is finite, as every finitely generated nilpotent torsion group is finite of bounded order (see, for instance,~\cite[Lemma~2.34]{CeSi}). Our proof of \cref{prop:bounds} does not rely on the primality of $\ell$. Hence, it can be generalized to arbitrary even girth $2\ell$ whether the following question has an affirmative answer.
\begin{question}
\label{question: general}
    Is \cref{cor: p-power-nilpotent} true for every positive integer $\ell$?
\end{question}

Last, we note that the bound from \cref{prop:bounds} is disproportionally high compared to the few examples we found in the census. Hence, we conclude by proposing the following problem.
\begin{question}
    What is the number of vertices of the smallest connected vertex-transitive $3$-valent and flexible graph of girth $2\ell$?
\end{question}


\subsection{Structure of the paper} \cref{sec:split,sec:BassSerre,sec:finiteIndex} contain the proof of \cref{thm} divided in various lemmas, while \cref{cor: p-power-nilpotent} and \cref{prop:bounds} are proved in \cref{sec: effective}.

We present here, in broad terms, the proof of \cref{thm}. The first step involves transforming the problem from the $3$-valent case to the $4$-valent case. This is achievable through the operation of splitting, which is explained in \cref{sec:split}. The advantage of the $4$-valent case lies in its easiness of handling with Bass--Serre theory. In fact, we can construct an infinite tree whose automorphism group is an amalgamated free product that encodes the structure of a $4$-valent graph where each vertex-stabilizer is isomorphic to the dihedral group of degree $4$. This is the subject of \cref{sec:BassSerre}. The task then becomes finding an appropriate normal subgroup $N$ of this amalgamated free product such that the $3$-valent split of the $4$-valent quotient graph possesses the properties prescribed by \cref{thm}. In rough terms, this means that almost all the elements in $N$ are long words with respect to the natural length of amalgamated products. The construction of this subgroup $N$ is addressed in \cref{sec:finiteIndex}.


\subsection{Acknowledgments}
We would like to thank Prof.~Ademir Hujdurović for posing us the question of whether there exist finite connected vertex-transitive $3$-valent and flexible graphs of arbitrary girth, for suggesting that the operation of splitting could aid in addressing this problem, and for sharing with us the preliminary version of~\cite{Ademir}. We would also like to thank Prof.~Montserrat Casals-Ruiz for helpful conversations.


\section{The operation of splitting}\label{sec:split}
    The operations of \emph{splitting} and \emph{merging} have been introduced in~\cite[Constructions~7 and~11]{PotocnikSpigaVerret2013} to create a framework for translating results from $3$-valent graphs with a perfect matching invariant under the action of their automorphism groups into results for $4$-valent graphs with a $2$-factor invariant under the action of their automorphism groups, and \textit{vice versa}. The fact that these operations are the inverse of one another (on both sides) has been proved in~\cite[Theorem~12]{PotocnikSpigaVerret2013} and~\cite[Theorem~2.9]{BarbieriGrazianSpiga2025}, and it suggests that we can reduce the proof of \cref{thm} to a problem on $4$-valent graphs with a $G$-invariant cycle decomposition. For our purposes, we only need to give a precise definition of the operation of splitting.

    \begin{construction}\label{con:splitting}
        Let $\Delta$ be a $4$-valent graph, and let $\mathcal{C}$ be a partition of $E\Delta$ into cycles. We build a $3$-valent graph, $\mathrm{s}(\Delta,\mathcal{C})$, whose vertex-set is
        \[ V\mathrm{s}(\Delta,\mathcal{C}) := \left\{(\alpha,\mathbf{c}) \in V\Delta \times \mathcal{C} \mid \alpha \in V\mathbf{c} \right\} \,\]
        and such that two vertices $(\alpha,\mathbf{c})$ and $(\beta,\mathbf{d})$ are declared adjacent if either $\mathbf{c}$ and $\mathbf{d}$ are distinct and $\alpha= \beta$, or $\mathbf{c}=\mathbf{d}$ and $\alpha$ and $\beta$ and are adjacent in $\mathbf{c}=\mathbf{d}$. The graph $\mathrm{s}(\Delta,\mathcal{C})$ is the \emph{split graph of the pair $(\Delta,\mathcal{C})$}.
    \end{construction}

    We need to point out some properties of \cref{con:splitting} that will be used in the following.
    
    Observe that, since $\Delta$ is $4$-valent, there are precisely two cycles in $\mathcal{C}$ passing through $\alpha$. Consequently, $\mathrm{s}(\Delta,\mathcal{C})$ is $3$-valent and
    \[ |V\mathrm{s}(\Delta,\mathcal{C})| = 2 |V\Delta| \,.\]
    
    Let $G\le \Aut(\Delta)$ be a group of automorphisms of $\Delta$. Note that, if $G$ fixes the edge-partition $\mathcal{C}$ setwise, then $G$ is a subgroup of $\Aut(\mathrm{s}(\Delta,\mathcal{C}))$. Moreover, if $G$ is arc-transitive on $\Delta$, then $G$ can swap the two cycles of $\mathcal{C}$ passing through each vertex. Hence, $G$ acts transitively on $V\mathrm{s}(\Delta,\mathcal{C})$. Observe that the $G$-invariance of $\mathcal{C}$ forces the action of any vertex-stabilizer on the neighbourhood of the fixed vertex to be imprimitive. We say that $G$ is \emph{locally-imprimitive}, for short.
    
    For any vertex $(\alpha,\mathbf{c}) \in V\mathrm{s}(\Delta,\mathcal{C})$, we observe that the vertex-stabilizer of $(\alpha,\mathbf{c})$ must contain the intersection of the vertex-stabilizer of $\alpha \in V\Delta$ and the setwise stabilizer of the cycle $\mathbf{c}$. In particular, whenever $G$ is arc-transitive on $\Delta$, the order of the vertex-stabilizers in $\mathrm{s}(\Delta,\mathcal{C})$ is half that of those in $\Delta$. Furthermore, by a connectedness argument, these groups are both $2$-groups (see, for instance,~\cite[Lemma~3.1]{Spiga2014}).

    We have collected all the data about \cref{con:splitting} we need to proceed.
    
   \begin{lemma}\label{lemma:girthSplit}
            Let $\Delta$ be a $4$-valent graph, and let $G$ be an arc-transitive locally-imprimitive group of automorphisms of $\Delta$ (or, equivalently, every vertex-stabilizer in $G$ is a $2$-group). Suppose that $\mathcal{C}$ is a $G$-invariant $2$-factor that contains girth cycles. Then
            \[ \g (\Delta) = \g \left( \mathrm{s}(\Delta, \mathcal{C})\right) \,. \]
        \end{lemma}
    \begin{proof}
    Consider the bijection $\theta: V\Delta \times \mathcal{C} \to V\mathrm{s}(\Delta, \mathcal{C})$ that \cref{con:splitting} induces. Note that, for every $\mathbf{c} \in \mathcal{C}$, $\theta(V\mathbf{c},\mathbf{c})$ is a cycle in $\mathrm{s}(\Delta,\mathcal{C})$ with the same length as $\mathbf{c}$. Hence, we immediately obtain the inequality
    \[ \mathrm{girth} (\Delta) \ge \mathrm{girth} \left( \mathrm{s}(\Delta, \mathcal{C})\right) \,. \]
    
    Aiming for a contradiction, we suppose that
    \[ \mathrm{girth} (\Delta) > \mathrm{girth} \left( \mathrm{s}(\Delta, \mathcal{C})\right) \,. \]
    Let $\mathbf{d}$ be a girth cycle in $\mathrm{s}(\Delta, \mathcal{C})$. Since all the cycles in $\mathcal{C}$ are longer than $\mathbf{d}$, it follows that the vertices
    \[ V\theta^{-1}(\mathbf{d}) := \{ \alpha \in V\Delta \mid \theta(\alpha,\mathbf{c}) \in V\mathbf{d}, \hbox{ for some }\mathbf{c}\in \mathcal{C} \}\]
    form a cycle of $\Delta$ whose length is strictly less than $\mathrm{girth}(\Delta)$, a contradiction. Hence, the equality is attained, and the proof is complete.
    \end{proof}

As a consequence of \cref{lemma:girthSplit}, we can focus our attention on locally-imprimitive $4$-valent graphs.
    
\section{The \texorpdfstring{$4$}{4}-valent scenario}\label{sec:BassSerre}
In this section, our aim is to build, for every integer $\ell \ge 2$, a $4$-valent graph whose automorphism group is arc-transitive and locally-imprimitive, and whose girth is $2\ell$. We will make use of Bass--Serre theory: we refer to~\cite{DicksDunwoodyGroups,SerreTrees} for the undefined terminology.

Let $G$ be the amalgamated product $D_4 *_{C_2} (C_2\times C_2)$, where
\[ D_4 = \langle a,b \mid a^4, b^2, \, a^b=a^{-1} \rangle \quad \hbox{and} \quad  C_2 \times C_2 = \langle ab \rangle \times \langle z \rangle \, \]
and the amalgamation is over $ \langle ab \rangle$. Recall that every $g\in G$ can be uniquely written in the \emph{reduced form}
\[ g = (ab)^\varepsilon a^{e_1}za^{e_2} \ldots a^{e_{i-1}}za^{e_i} \,,\]
where $\varepsilon\in \{0,1\}$, $i \in \mathbb{N}$, $e_1,  e_i\in \{0,1,2,3\}$, and $e_2, \dots , e_{i-1}\in \{1,2,3\}$
(see, for instance,~\cite[Theorem~I.1]{SerreTrees}).
Following~\cite{SerreTrees}, we define the \emph{length of $g$} to be the number of symbols $a$, $a^2$, $a^3$ and $z$ we use to write it in its reduced form. Observe that this length defines a pseudometric on $G$ in the usual fashion, and hence we can talk about balls of centre $1$,
\[ \B_\ast(k) := \left\{ g \in G \mid \len_\ast(g)\le k \right\} \,.\]

We now state a result that we will prove in \cref{sec:finiteIndex}.
\begin{lemma}\label{lemma:finiteIndexNormal}
    Let $G$ be the amalgamated product $D_4 *_{C_2} (C_2\times C_2)$, where
    \[ D_4 = \langle a,b \mid a^4, b^2, \, a^b=a^{-1} \rangle \quad \hbox{and} \quad 
    C_2 \times C_2 = \langle ab \rangle \times \langle z \rangle \, \]
    and the amalgamation is over $ \langle ab \rangle$. Then, for every prime $\ell$, there exists a finite index normal subgroup $N$ of $G$ such that
    \begin{enumerate}[$(a)$]
        \item $N \cap \B_\ast(4\ell -1) = \left\{1\right\}$,
        \item $N \cap \B_\ast(4\ell + 1) = \left\{1, \, (za^2)^{\pm 2\ell} \right\}$.
    \end{enumerate}
\end{lemma}

We now start the construction of our putative $4$-valent graph $\Delta$. We consider the graph of groups
    \[\begin{tikzcd}
        D_4 &&& C_2 \times C_2
        \arrow["C_2", no head, from=1-1, to=1-4]
    \end{tikzcd}\]
and we denote by $\mathcal{T}$ the standard tree obtained from this graph of groups. Recall that $\mathrm{Aut}(\mathcal{T})$ is isomorphic to $G$, and that, as $|D_4:C_2|=4$ and $|C_2\times C_2:C_2|=2$, the vertices labelled by right cosets of $D_4$ have valency $4$, while those labelled by right cosets of $C_2\times C_2$ have valency $2$.

Observe that there is a one-to-one correspondence between non-backtracking paths of even length starting in $D_4$ and right cosets of $D_4$. Moreover, the length of one of these paths coincides with the (minimal) length of the reduced form among the representatives of the right coset of $D_4$ which is the end-vertex of the path. 

Note that $\mathcal{T}$ is the barycentric subdivision of the infinite $4$-valent tree $\mathcal{S}$ (that is, by exchanging each vertex labelled by right cosets of $C_2\times C_2$ with an edge joining its two neighbours, we obtain the infinite $4$-valent tree). Therefore, $G$ acts arc-transitively on $\mathcal{S}$.
The metrics of these graphs are such that two vertices at distance $2h$ in $\mathcal{T}$ are at distance $h$ in $\mathcal{S}$.

Before proceeding, it is convenient to have a concrete example of a path in $\mathcal{S}$ spelled out. For instance, the word $za^{e_1}za^{e_2}$ defines the path
\[\begin{tikzcd}
    D_4 && D_4 za^{e_2} && D_4 za^{e_1}za^{e_2}
    \arrow[no head, from=1-1, to=1-3]
	\arrow[no head, from=1-3, to=1-5]
\end{tikzcd}\]
while its inverse $a^{-e_2}za^{-e_1}z$ corresponds to the picture
\[\begin{tikzcd}
    D_4 && D_4 z && D_4 za^{-e_1}z
    \arrow[no head, from=1-1, to=1-3]
	\arrow[no head, from=1-3, to=1-5]
\end{tikzcd}\]

Let $N$ be the finite index normal subgroup of $G$ described in \cref{lemma:finiteIndexNormal}, and let $\Delta_\ell = \mathcal{S}/N$ be the finite $4$-valent graph obtained by quotienting $\mathcal{S}$ by $N$ (that is, the graph whose vertex-sets and edge-sets coincides with the $N$-orbits of $V\mathcal{S}$ and $E\mathcal{S}$, respectively). We want to show that $\Delta_\ell$ satisfies the following additional properties.

\begin{lemma}\label{lemma:BassSerre}
    For every positive integer $\ell\ge 2$, there exists a finite connected $4$-valent graph $\Delta_\ell$ such that
    \begin{enumerate}[$(a)$]
        \item $\Aut(\Delta_\ell)$ is arc-transitive and locally-imprimitive,
        \item each vertex-stabilizer has order at least $8$,
        \item the $\mathrm{Aut}(\Delta_\ell)$-invariant $2$-factor consists of girth cycles,
        \item $\g(\Delta_\ell)=2 \ell$.
    \end{enumerate}
\end{lemma}
\begin{proof}
Let $\alpha$ be the vertex of $\Delta_\ell$ where the vertex $D_4$ of $\mathcal{T}$ projects to through the operations we have previously described. First, we observe that $\mathrm{Aut}(\mathcal{T})/N = G/N$ is isomorphic to a subgroup of $\mathrm{Aut}(\Delta_\ell)$. Observe that $D_4 \subseteq \B_\ast(1)$, thus \cref{lemma:finiteIndexNormal}~$(a)$ implies that $D_4 \cap N$ is trivial. It follows that the vertex-stabilizer of $\alpha$ contains
\[G_{D_4}/N \cong D_4N/N \cong D_4 \,,\]
which is a subgroup of order $8$. This proves $(b)$.
    
Let us denote by $\alpha_z$, $\alpha_{za}$, $\alpha_{za^2}$ and $\alpha_{za^3}$ the neighbours of $\alpha$, which can be recognized as the projection of the vertices $D_4 z$, $D_4 za$, $D_4 za^2$ and $D_4 za^3$, respectively. Observe that the action of $D_4$ on the neighbourhood of $\alpha$ defines the blocks of imprimitivity $\{\alpha_z, \alpha_{za^2}\}$ and $\{\alpha_{za}, \alpha_{za^3}\}$. Furthermore, using \cref{lemma:finiteIndexNormal}~$(b)$ and recalling the correspondence between non-backtracking paths and words, the girth cycles through $\alpha$ in $\Delta_\ell$ have length $2\ell$ and they contain either the pair $\{\alpha_z, \alpha_{za^2}\}$ or $\{\alpha_{za}, \alpha_{za^3}\}$ (the former being induced by the word $(za^2)^{\pm 2\ell}$ and the latter by $(a^{-1}(za^2)^{2\ell} a)^{\pm 1}$). Therefore, $(d)$ holds.

Moreover, \cref{lemma:finiteIndexNormal}~$(b)$ also implies that there exists no girth cycle that contains any pair of neighbours of $\alpha$ distinct from $\{\alpha_z, \alpha_{za^2}\}$ and $\{\alpha_{za}, \alpha_{za^3}\}$. It follows that $\mathrm{Aut}(\Delta_\ell)_\alpha$ must define the same two block of imprimitivity on the neighbourhood of $\alpha$ as $D_4$ does, and hence $(a)$ holds. Last, we can note that the $\mathrm{Aut}(\Delta_\ell)$-invariant $2$-factor contains precisely the $\mathrm{Aut}(\Delta_\ell)$-orbits of these two girth cycles through $\alpha$. Thus, $(c)$ is proved.
\end{proof}

We are now fully equipped to prove the main theorem of this note.

\begin{proof}[Proof of \cref{thm}]
    For every $\ell\ge 2$, let $\Delta_\ell$ be the $4$-valent graph described in \cref{lemma:BassSerre}, and let $\mathcal{C}$ be the $\Aut(\Delta_\ell)$-invariant $2$-factor of $\Delta_\ell$. Using \cref{lemma:girthSplit} and the fact that the order of every vertex-stabilizer is halved by applying \cref{con:splitting}, we obtain that
    \[ \Gamma_\ell = \mathrm{s}(\Delta_\ell, \mathcal{C})\]
    enjoys all the properties required by \cref{thm}.
\end{proof}


\section{Subgroups consisting of long words}\label{sec:finiteIndex}
\cref{sec:finiteIndex} is devoted to proving \cref{lemma:finiteIndexNormal}. We will need some preliminary results.

Let $G$ be a finitely generated group on generators $S = \{ x_1, \dots, x_k\}$, and let $g\in G$. The \emph{length of \( g\) with respect to $S$} is the length of the shortest word in the elements of \( S \) that represents \( g \), or, equivalently, the distance between $1$ and $g$ in the Cayley graph $\mathrm{Cay}(G,S)$. We denote the length of $g$ by \(\len_S(g)\). In the same fashion, we will denote by $\B_S(\ell)$ the ball of radius $\ell$ around the identity (usually when $G$ is the free group, we will drop the subindexes $S$, as it will refer to the canonical generators). Note that this notation is consistent with that introduced in Section~\ref{sec:BassSerre}, because, in that context, it corresponds to the length in the quotient $(D_4 *_{C_2} C_2 \times C_2) / C_2$ with respect to the generating set $\{a, a^2, a^3, z \}$.

\begin{lemma}
\label{lem: res fin}
Let $G$ be a finitely generated group with finite generating set $S$, and let $H \unlhd G$ be a normal subgroup. Suppose that $G/H$ is residually finite. Then, for every positive integer $\ell$, there exists a finite index normal subgroup \( K_\ell \trianglelefteq G \) containing \( H \) such that 
\[ K_\ell \cap \B_S(\ell) = H \cap \B_S(\ell).
\]
\end{lemma}
\begin{proof}
Since $G/H$ is residually finite, there exists a family of finite index normal subgroups $\{ H_i \}_{i \in I}$ such that 
\[ H = \bigcap_{i \in I} H_i.\]
For each $w \in \B_S(\ell) - H$, there exists $i(w) \in I$ such that $w$ is not an element of $H_{i(w)}$. Accordingly, consider the finite index normal subgroup 
\[ K_\ell = \bigcap_{w \in \B_{S}(\ell) - H} H_{i(w)}.\]
Then \( H \leq K_\ell \) and $K_\ell$ and $\B_S(\ell) - H$ are disjoint. Therefore, $K_\ell$ is the required normal subgroup.
\end{proof}

We require a small-scale form of \Cref{thm: p-power} for arbitrary integers.

\begin{lemma}
\label{lem: powers 1}
Let $F_3 = \langle x_1, x_2, x_3 \rangle$ be the free group of rank $3$, and let $H$ be the normal closure of the subgroup generated by $x_1^\ell$ and $(x_2x_3^{-1})^\ell$, that is, $H = \langle x_1^\ell, (x_2x_3^{-1})^\ell \rangle^{F_3}$. Then every word $w$ in $H$ satisfies $\len(w) \geq \ell$, and equality is attained if and only if $w = x_1^{\pm \ell}$.
\end{lemma}
\begin{proof}
Consider the natural quotient 
\[ F_3/H = \left\langle y_1, \, y_2, \, y_3 \,\middle\vert\, y_1^{\ell}, \, \left(y_2y_3^{-1} \right)^\ell \right\rangle \,.\]
We rewrite its generators by putting
\[z_1 = y_1 , \quad z_2=y_2y_3^{-1} , \quad z_3=y_3 \,.\]
Hence, we see that
\[ F_3/H = \left\langle  z_1, \, z_2, \, z_3 \,\middle\vert\, z_1^\ell, \, z_2^\ell  \right\rangle = C_\ell * C_\ell * C_\infty \,\]
is the free product of three cyclic groups. 
Consider the projection map $\pi: F_3 \to F_3/H$ defined on the generators by
\[x_i ^\pi = z_i, \quad \hbox{for every } i\in\{1,2,3\} \,.\]
Let $w = x_{i_1}^{\alpha_1} \dots x_{i_k}^{\alpha_k}$ be a nontrivial reduced word, for some exponents $\alpha_j \in \Z$ and indices $i_j \in \{1, 2, 3 \}$ so that adjacent letters are distinct.

Assume that $\len(w) < \ell$. Then,
\[ w^\pi = z_{i_1}^{\alpha_{1}} \dots z_{i_k}^{\alpha_k} \,.\]
As $\len(w) < \ell$, each $\alpha_j \not\equiv 0 \pmod{\ell}$, and, hence, $w^\pi$ is nontrivial. Since $H=\ker(\pi)$, $w$ is not an element of $H$, as required.

Assume now that $\len(w)=\ell$ and that $w \in H$. As $w^\pi$ is trivial and $w$ is not, we have that $w$ projects to either $z_1^{\pm\ell}$ or $z_2^{\pm\ell}$. In the former case, $w = x_1^{\pm \ell}$. In the latter,  $w = (x_2x_3^{-1})^{\pm \ell}$, against the hypothesis on the length. Therefore, $w=x_1^{\pm\ell}$, which completes the proof.
\end{proof}

The last piece of information we need for our proof is the following partial result  concerning Baumslag's Conjecture on the residual finiteness of one-relator groups.
\begin{thm}[\cite{AT}, Theorem~2]
\label{thm: AT}
    Let
    \[ G = \left\langle \, x, y \mid (x^\alpha y^\beta x^\gamma y^\delta)^\ell \, \right\rangle \]
    where $\alpha, \beta, \gamma, \delta$ are integers, and $\ell$ is a positive integer. Then, $G$ is residually finite.
\end{thm}

We are now ready to dive into the bulk of \cref{sec:finiteIndex}.
\begin{proof}[Proof of \cref{lemma:finiteIndexNormal}]
    We start by building a finite index normal free subgroup\footnote{We point out that every amalgamated free product whose factors are finite contains a finite index free normal subgroup that can be explicitly constructed: see, for instance,~\cite[Theorem~I.7.4, Proposition~I.7.10, Theorem~IV.1.6 and Corollary~IV.1.9]{DicksDunwoodyGroups}.} of $G$.  We consider the group homomorphism $\varphi: G \to \mathrm{Sym}(8)$
    defined on the generators by
    \[ a^\varphi = (1 \,2 \, 3 \, 4)(5 \,8 \,7 \,6), \quad b^\varphi = (1\,2)(3\,4)(5\,6)(7\,8), \quad z^\varphi = (1\,2)(3\,4)(5\,8)(6\,7) \,.\]
    Observe that the restrictions of $\varphi$ to $D_4 = \langle a,b\rangle$ and to $C_2\times C_2 = \langle ab, z\rangle$ are isomorphisms. Therefore,~\cite[Proposition~I.7.10]{DicksDunwoodyGroups} implies that $\ker(\varphi)$ is a free group. By the first isomorphic theorem, its index in $G$ is $16$, and thus, by~\cite[Remark~IV.1.11]{DicksDunwoodyGroups}, the rank of $\ker(\varphi)$ is $3$. To find a minimal set generators of $\ker(\varphi)$ we implement $\varphi$ in \textsf{GAP}, and we computed
    \[ \ker(\varphi) = \langle  za^2za^2, zaza, za^3za^3 \rangle \,.\]
    From here on, we identify $\ker(\varphi)$ with $F_3 = \langle x_1,x_2,x_3 \rangle$ via the correspondence
    \[ x_1 \mapsto za^2za^2, \quad x_2 \mapsto zaza, \quad x_3 \mapsto za^3za^3 \,, \]
    and we set $S = \{x_1,x_2,x_3\}$.
    
    Let $H = \langle x_{1}^\ell\rangle^G \le F_3$ be the normal closure of $ \langle x_{1}^\ell\rangle$ in $G$. Note that $G$ embeds in the automorphism group of $F_3$. As a consequence, the conjugacy classes of $x_1^{\ell}=(za^2)^{ 2\ell}$ in $G$ can be described as the orbit of $(za^2)^{2\ell}$ under the action of a prescribed subgroup of Nielsen transformations. Elementary computations show that
    \begin{gather*}
        x_1^a = x_2^{-1}x_3 \,, \quad (x_2^{-1}x_3)^a = x_1^{-1} \,, \quad x_1^b = x_3^{-1}x_2 \,,
        \\
        (x_2^{-1}x_3)^b = x_1 \,, \quad x_1^z=x_1^{-1}\,, \quad (x_2^{-1}x_3)^z = x_3^{-1}x_2 \,.
    \end{gather*}
    This facts can be summarized in the following diagram.
    \[\begin{tikzcd}
	   && {x_1} \\
	   \\
	   {x_2^{-1}x_3} &&&& {x_3^{-1}x_2} \\
	   \\
	   && {x_1^{-1}}
	   \arrow["a"', curve={height=-12pt}, from=1-3, to=3-1]
	   \arrow["b", curve={height=-12pt}, tail reversed, from=1-3, to=3-5]
	   \arrow["z", curve={height=-6pt}, tail reversed, from=3-1, to=3-5]
    	\arrow["a"', curve={height=-12pt}, from=3-1, to=5-3]
    	\arrow["b"', curve={height=18pt}, tail reversed, from=3-1, to=5-3]
    	\arrow["a"', curve={height=-12pt}, from=3-5, to=1-3]
    	\arrow["z"', curve={height=6pt}, tail reversed, from=5-3, to=1-3]
    	\arrow["a"', curve={height=-12pt}, from=5-3, to=3-5]
    \end{tikzcd}\]
    In particular, we can conclude that
    \[ \left( x_1^{\ell} \right)^G = \left\{ x_1^{\pm \ell},\,(x_2^{-1}x_3)^{\pm \ell} \right\} ^{F_3} \,\]
    and thus
    \[ H = \langle x_1^{\pm \ell} \rangle ^G = \langle x_1^{\ell},\, (x_2^{-1}x_3)^{\pm \ell} \rangle ^{F_3} \,.\]

   In order to apply~\Cref{lem: res fin}, we shall prove that $G/ H$ is residually finite. For that purpose, observe that $G/H$ is a finite extension of 
    \[  F_3 / H = \langle x_1, \, x_2, \, x_3 \mid x_1^\ell, \, (x_2 x_3^{-1})^\ell  \rangle \cong C_\ell \ast B_\ell  \,,\]
    where $C_\ell$ is a cyclic group of order $\ell$, and $B_\ell = \langle a, b \mid (ab^{-1})^\ell \rangle$. Hence, it is enough to prove that $C_\ell \ast B_\ell$ is itself residually finite.\footnote{Let $G$ be a virtually residually finite group, and let $H$ a finite index normal subgroup of $G$ which is residually finite. If $g \in G - H$, then $g$ has a nontrivial image in the canonical epimorphism $G \to G/H$. If $g \in H$, let $\varphi : H \to F$ be an homomorphism such that $g$ is not in $\ker(\varphi)$, and let $K$ be the normal core of $\ker(\varphi)$ in $G$. Observe that
    \[|G:K| \le |G: \ker(\varphi)| !. \]
    and hence it is finite. In particular, $G/K$ is a finite group. It follows that $g$ has a nontrivial image in the canonical epimorphism $G \to G/K$. This proves that a virtually residually finite group is residually finite, and, as every free group is residually finite, that a virtually free group is residually finite.} Indeed, $C_\ell$ is finite, $G_\ell$ is residually finite by Theorem~\ref{thm: AT}, and the free product of residually finite groups is also residually finite. Therefore, \cref{lem: res fin} implies the existence of a finite index normal subgroup $K_\ell \leq G$ such that 
    \[ H \cap \B_S(\ell) = K_\ell \cap \B_S(\ell). \]

    To conclude the proof, we need to verify how $K_\ell$ intersects with the balls $\B_\ast (k)$ of $G$.  We start by establishing the elements contained in $K_\ell \cap \B_S(\ell)$, and to do so we leverage on the previous identity. By \Cref{lem: powers 1}, for every element $w \in H$, $\len(w) \geq \ell$, and, if $\len(w) = \ell$, then $w = x_1^{\pm \ell}$.  Accordingly, we conclude that
    \[ \B_S(\ell) \cap K_\ell = \B_S(\ell) \cap H = \{1,\, x_1^{\pm \ell} \} \,.\]
    It follows at once that
    \[ \B_\ast(4\ell-1) \cap K_\ell = \{ 1 \} \quad \hbox{and} \quad \B_\ast(4\ell) \cap K_\ell = \{ 1, \, (za^2)^{\pm 2\ell} \} \,.\]
    Finally, by eventually substituting $K_\ell$ for $K_\ell \cap F_3$, we observe that every element in $K_\ell$ contains an even number of symbols $z$, and thus
    \[ \B_\ast(4\ell) \cap K_\ell = \B_\ast(4\ell+1) \cap K_\ell = \{ 1, \, (za^2)^{\pm 2\ell} \} \,.\]
    Hence, $K_\ell$ enjoys all the properties required. This concludes the proof. 
\end{proof}

\section{Effective bound for primes}
\label{sec: effective}
To obtain an effective bound on the number of vertices of the graphs $\Gamma_\ell$, it is necessary to control the index of the normal subgroup of prescribed girth appearing in \Cref{lemma:finiteIndexNormal}. Since the problem reduces to free groups, we study the girth of certain verbal subgroup of the free group.

We start with the girth of the $\ell$-power subgroup
\[ F_k^\ell = \langle g^\ell \mid g \in F_k \rangle.\]
The following result seems to be of   elementary nature within geometric group theory, yet we have not been able to locate a reference in the existing literature. We include a proof in the case where $\ell$ is prime. Furthermore, asymptotically for large $\ell$, the same result can be derived by standard arguments from Small Cancellation Theory.

\begin{proposition} \label{thm: p-power}
    Let $F_k = \langle x_1, \dots, x_k \rangle $ be the free group of rank $k$, and let $\ell$ be a prime number. 
    Then, $\len(w) \geq \ell$ for every $w \in F_k^\ell$, and equality is attained if, and only if, $w = x_{i}^{\pm \ell}$, for some index $i \in \{ 1, \dots, k \}$.
\end{proposition}
\begin{proof}
We aim to prove that there exists a $k$-generated group $G$ of exponent $\ell$ such that any nontrivial element $w \in \B(\ell-1)$ is not a law. It follows that there exists a suitable normal subgroup $N \unlhd F_k$ such that \( F_k / N \cong G\), $F_k^\ell \leq N$ and $\B(\ell -1) \cap N = \{1\}$.

Let $w = x_{i_1}^{\alpha_1} \dots x_{i_n}^{\alpha_n}$ be a nontrivial reduced word, for some $i_j \in \{ 1, \dots, k \}$ and $\alpha_j \in \Z$. Assume that $\len(w) < \ell$, so that $|\alpha_i| < \ell$ for all $i$. Consider the Magnus embedding
\[
\mathcal{M}\colon F_k\to \mathbb{Z}\langle\!\langle X_1,\dots,X_k\rangle\!\rangle,
\qquad  x_i \mapsto 1+X_i,
\]
into the $\mathbb{Z}$-algebra of noncommutative formal power series in $k$ variables (see~\cite[Proposition 10.1]{LyndonSchuppCombinatorial}). As in~\cite[Section~2.2]{SuWa}, the coefficient $c_W$ of the monomial $W = X_{i_1} \dots X_{i_n}$ in $\M(w) $ is given by 
\[ \varepsilon(\partial_{i_1} \dots \partial_{i_n}(w)) \,, \]
where $\partial_j$ is the Fox derivative $\partial_j \colon \Z[F_k] \rightarrow \Z[F_k]$ defined as $\partial_j(x_i) = \delta_{ij}$ and $\varepsilon \colon \mathbb{Z}[F_k] \rightarrow \Z$ is the augmentation map (see~\cite[Chapter I.10]{LyndonSchuppCombinatorial}). Moreover,~\cite[Equation (4)]{MaPu} implies
\[ c_W = \varepsilon(\partial_{i_1} \dots \partial_{i_n}(w)) = \prod_{i=1}^n \alpha_i \,.\]

Let $\F_\ell$ be the finite field of $\ell$ elements, let
\( \M_\ell \colon F_k \rightarrow \F_\ell \langle \langle X_1,\dots,X_k \rangle \rangle\) be the Magnus embedding where the coefficients are reduced modulo $\ell$,
and let 
\[ \mathcal{I}_\ell = \left\{ f \in \F_\ell \langle \langle X_1,\dots,X_k \rangle \rangle \mid \mathrm{ord}(f) \geq \ell  \right\} \footnote{Here, $\mathrm{ord}(f)$ is the \emph{order} of $f$, {\it i.e.}, the lowest degree of a monomial appearing in the power series $f$.}\]
({\it cf.}~\cite[Section 3.1]{Zozaya}).
We observe that
\[\deg(W) \leq \len(w) < \ell \,,\]
and, as $|\alpha_i| < \ell$ for every $i\in\{1,\dots,k\}$, $c_W$ and $\ell$ are coprime.
We obtain that \[\M_\ell(w) \notin \mathcal{I}_\ell \,.\]

We now define $G = \M_\ell(F_k) / \mathcal{I}_\ell$.
To complete the proof, we will show that $G$ is a group of exponent $\ell$. For every $g \in F_k$,
\[\M_\ell(g^\ell) = \M_\ell(g)^\ell = (1 + f)^\ell = 1+ f^\ell \,,\]
for some power series $f$ such that $f(0, \dots, 0) = 0$. Write $f$ as the sum of monomials $\sum f_i$, where $\deg f_i \geq 1$, then
\[\deg(f_1 \dots f_{\ell}) \geq \sum_{i=1}^{\ell} \deg(f_i) \geq \ell \,.\]
Thus, $f^\ell \in \mathcal{I}_\ell$. This proves that $G$ has exponent $\ell$. 

We now characterize equality. Suppose that $w \in F_k^\ell$ satisfies $\len(w) = \ell$. If $\ell$ is odd, then $w \notin \gamma_2(F_k) = [F_k,F_k]$. Consider the abelianisation map $\pi \colon F_k \to \mathbb{Z}^k$.  
If $w \in F_k^\ell - \gamma_2(F_k)$, then $w^\pi \in \ell \mathbb{Z}^k - \{0 \}$, and hence
\[
    \ell = \len (w) \geq \len_{\{x_1^\pi, \dots,x_k^\pi\}}(w^\pi) \geq \ell.
\]
Therefore, $w^\ell = \pm \ell \,x_i^\pi$ for some $i \in \{1, \dots, k\}$. Hence, $w = x_i^{\pm \ell}$. Finally, if $\ell = 2$, there exists no element $w\in F_k$ with $\len(w) = 2$ that lies in $\gamma_2(F_k)$, and thus the same argument applies.  
\end{proof}

We are ready to prove \cref{cor: p-power-nilpotent}. Moreover, we remark that this result cannot be obtained from elementary considerations of Small Cancellation Theory.
\begin{proof}[Proof of \cref{cor: p-power-nilpotent}]
Keeping the notation of the preceding proof, for every $w \in \gamma_\ell(F_k)$, we have $\M(w) \in \mathcal{I}_\ell$ (see~\cite[Proposition~10.2]{LyndonSchuppCombinatorial}). Hence, the group $G$ is $\ell$-step nilpotent, and thus there exists a normal subgroup $N \unlhd F_k$ such that $F_k^\ell \gamma_\ell(F_k) \leq N$ and $\B(\ell-1) \cap N=\{1 \}$.

To establish equality, we observe again that if $w \in F_k^{\ell} \gamma_\ell(F_k)$ and $\len(w) = \ell$, then $w \notin \gamma_2(F_k)$. From here, we can conclude as in the previous proof. 
\end{proof}

\begin{proof}[Proof of \Cref{prop:bounds}]
Following the proof of \Cref{lemma:finiteIndexNormal}, let $F_3 \unlhd G$ be the finite index normal free subgroup of rank $3$ such that $|G: F_3| = 16$, and let 
\[  x_1 = za^2za^2, \, \quad x_2 = zaza, \, \quad x_3 = za^3za^3 \]
be the canonical generators for $F_3$.
Recall from the same proof that 
we have that
\[ \left( x_1^{\ell} \right)^G = \left\{ x_1^{\pm \ell},\,(x_2^{-1}x_3)^{\pm \ell} \right\} ^{F_3}. \]
From here, the proofs diverge, as here we are building an explicit finite index subgroup $N \leq F$ that is normal in $G$, and such that
\[\B_S(\ell) \cap N = \left\{ 1, x_{1}^{\pm \ell} \right\} \,.\]

Define $N=F_3^{2 \ell} \gamma_{\ell}(F_3) \leq F_3^\ell \gamma_\ell(F_3)$. Observe that, as $N$ is a verbal subgroup of $F$, it is characteristic in $F$, and hence normal in $G$. Let
\[ M = \langle x_1^\ell, N\rangle^G = \langle x_1^\ell, (x_2 x_3^{-1})^\ell, N \rangle  \leq F_3^{\ell}\gamma_{\ell}(F_3).\]
On the one hand, \Cref{cor: p-power-nilpotent} states that $\len(w) \geq \ell$ for every $w \in M$. On the other hand, if $\len(w) = \ell$, then $w=x_i^{\pm \ell}$ for some $i \in \{1,2,3\}$. However, we claim that $x_2^\ell$ and $x_3^\ell$ are not elements of $M$. Aiming for a contradiction, suppose that $x_2^\ell \in M$, that is,
\begin{equation} \label{eq:ekis2}
    x_2^\ell = x_1^{g_1 \ell} (x_2 x_{3}^{-1})^{h_1 \ell} \dots x_1^{g_m \ell} (x_2 x_3^{-1})^{h_m \ell} u v \,,
\end{equation}
with $m \geq 1$, $g_i, h_i \in F_3$, $u \in F_3^{2 \ell}$ and $v \in \gamma_\ell(F_3)$.
Consider the evaluation epimorphism $\varphi: F_3 \rightarrow \mathbb{Z}$ defined on the generators by
\[ x_1 \mapsto 1 \,, x_2 \mapsto x_2, \, x_3 \mapsto x_2 \,.\]
Upon applying $\varphi$ to \cref{eq:ekis2}, we obtain
\[ x_2^\ell = (x_2^\ell)^\varphi = \left(x_1^{g_1 \ell} (x_2 x_{3}^{-1})^{h_1 \ell} \dots x_1^{g_m \ell} (x_2 x_3^{-1})^{h_m \ell} \right)^ \varphi u^\varphi v^\varphi = u^\varphi \in \left\langle x_2^{2\ell} \right \rangle \,,\]
a contradiction. Since $x_2^\ell$ is not an element of $N^\varphi$, the same is true for $N$ itself. The fact that $x_3^\ell$ is not an element of $N$ can be proved in a similar fashion. 

To conclude, we observe that
\[|V\Gamma_\ell| = |G : M| \leq |G: N| = 16 |F_3 : F_3^{2\ell} \gamma_\ell(F_3)| = 16 \cdot N(3,2\ell) \,. \qedhere\]
\end{proof}

\bibliographystyle{plain}
\bibliography{bibFlexible}

\end{document}

